\section{Conclusion}\label{sec:conclusion}

Mean--variance portfolio choice conventionally begins by estimating a return
covariance matrix, precisely the object that short and high-dimensional panels
make difficult to estimate reliably.
This paper takes a different route.
It asks what observed differences between firms' information distributions can
certify about portfolio risk before cross-asset covariance is estimated.
The resulting information certificate is a one-sided restriction on admissible
portfolio variance, not a point estimate of the covariance matrix.

Related work establishes neighbouring implications of the same distributional
geometry at other levels of aggregation.
At the pairwise level, \citet{gawronsky_continuous_2026} derive a covariance
envelope from Wasserstein separation.
At the cross-sectional level, \citet{gawronsky_spatial_2027} form a barycentric
interaction field through target-anchored Wasserstein barycentric
reconstruction.
The distinct step here is portfolio aggregation.
An antilipschitz carrier with bounded firm-specific slack transfers observed W2
separation into lower bounds on latent exposure separation; one coherent joint
exposure law then aggregates those restrictions into a portfolio-risk bound.
A maintained return bridge translates the exposure result into standardized and
raw-return variance statements.
The text determines the observed geometry, but it does not identify any of these
transmission restrictions.

Within this bridge, multi-firm transport dispersion supplies the sharp
certificate, while the weighted pairwise certificate provides the
computationally simpler decision rule used in the empirical analysis.
The investor minimizes the certified upper bound rather than an estimated
portfolio variance.
When firm-specific slack is zero, the common carrier scale changes how much
variance is certified away but not the selected normalized risk allocation,
which depends only on the observed W2 geometry.
Capital implementation still requires marginal volatility scales, but expected
returns and cross-asset return covariance do not enter the allocation's
construction.

In the \ppnum[0]{n-tickers}-firm 2018--2022 panel, the
Qwen3-Embedding-8B news-only allocation
falls between the
\ppnum[2]{news-only-percentile-uniform-cap-15-pct}\% and
\ppnum[2]{news-only-percentile-concentrated-cap-15-pct}\% variance percentiles
across four prespecified capped reference populations, while equal risk weights
rank higher under each corresponding comparison.
Its standardized variance is
\ppnum[1]{news-only-equal-variance-reduction-pct}\% below the variance of the
equal-risk benchmark but
\ppnum[1]{news-only-gmv-excess-variance-pct}\% above the ex post long-only
sample GMV.
The evidence is therefore an unusually low location among ordinary feasible
allocations, not covariance-free replication of the covariance-informed
optimum.
Portfolio anatomy further shows that
\ppnum[1]{anatomy-cross-sector-credit-pct}\% of the additive certificate credit
comes from cross-sector pairs.
This attribution locates the geometry used by the allocation; it does not
identify a sector mechanism.
Both the information geometry and the covariance used for evaluation draw on the
same 2018--2022 period.
The resulting ranking is therefore descriptive and in-sample: it illustrates
the allocation implied by the maintained framework rather than identifying the
information-to-risk transmission or predicting future portfolio variance.

As a theoretical restriction, the certificate can complement conventional
covariance estimation with information about admissible portfolio risk that
originates outside joint returns.
The present paper does not incorporate the certificate into a shrinkage or
constrained covariance estimator, nor does it establish the properties of such
an estimator.
Doing so is a natural extension when return histories are short or unstable.

More broadly, the construction is not intrinsically tied to corporate news or
equities.
Where economically relevant objects admit distribution-valued representations
and a defensible transmission restriction links their observed geometry to
latent risk, the same logic can generate one-sided restrictions on aggregate
risk.
The immediate empirical priorities are point-in-time construction, genuinely
out-of-sample evaluation, and identified or calibrated counterparts to the
maintained carrier, firm-specific slack, and return bridge.
Those extensions would test the certificate's practical value; the present
contribution is the coherent portfolio-level object and the explicit bridge on
which it depends.
